\subsection{Validación y aceptación de la admisión documental general}\label{sec:pruebas-admision-documental}

Las comprobaciones focales de la admisión general distinguen la política de
aplicación, el rechazo HTTP y la ejecución de analizadores nativos. La
\cref{tab:pruebas-admision-documental} presenta los resultados aprobados por
frontera; no representa una nueva ejecución global del proyecto ni actualiza
las mediciones históricas de cobertura. Los 29 casos de infraestructura y los
cinco del adaptador nativo son pruebas distintas, sin acumular sus repeticiones
focales como casos adicionales.

\begin{table}[H]
  \centering
  \caption{Pruebas focales de admisión de nuevas cargas y versiones.}
  \label{tab:pruebas-admision-documental}
  \begin{tabularx}{\linewidth}{@{}L{3.2cm} r X@{}}
    \toprule
    \textbf{Frontera} & \textbf{Aprobadas} & \textbf{Comprobación principal} \\
    \midrule
    Aplicación & 8 & Tres ingresos documentales, rechazo previo a persistencia y reautenticación. \\
    Cliente Node & 12 & Errores tipados, contexto de sesión y envío único. \\
    Navegador controlado & 8 & Rechazo de carga o versión y conservación de archivo y borrador. \\
    Adaptador nativo & 5 & Familias reales y corrupción interna de contenido comprimido. \\
    Proceso privado & 5 & Límites de recursos, descriptores heredados y protocolo. \\
    Infraestructura & 29 & Estructuras de archivo, inventario multimedia y supervisión de procesos. \\
    HTTP & $1+1$ & Una prueba unitaria y una de rutas: categorías públicas y transporte de errores. \\
    Instalador & 15 & Fuente, configuración y provisión comprobadas en VPS3. \\
    \bottomrule
  \end{tabularx}
\end{table}

Las pruebas nativas comprueban que no basta una envoltura válida cuando el
contenido comprimido interno está dañado. Las del proceso privado observan sus
límites, la retirada de descriptores ajenos y el rechazo de respuestas de
protocolo incompletas o inesperadas. El supervisor diferencia exceso de salida,
vencimiento, rechazo del archivo e indisponibilidad de ejecución. Los recorridos
con HTTP controlado comprueban que el fallo no borra el archivo seleccionado ni
los metadatos y que no provoca una segunda escritura automática. Esos recorridos
no sustituyen el ingreso contra PostgreSQL, Redis y el decodificador compuesto
en el servidor real.

La aceptación integrada mediante \rutarepo{scripts/api-demo.sh} terminó con
salida cero. Después del respaldo y la restauración, las ocho familias
conservaron sus bytes originales exactos. El navegador con servicios reales
aprobó cuatro escenarios en 27.2~s: uno nuevo de admisión documental, cuya
ejecución duró 8.4~s, y tres regresiones de contenido. Se inspeccionaron las
capturas de escritorio a 1440 píxeles y móvil a 390 píxeles, sin desbordamiento
horizontal. Estas evidencias complementan las focales; no acreditan una
campaña de usabilidad con personas.

La dependencia se verificó con el usuario del runner en los tres VPS. VPS1,
con AlmaLinux, requirió una compilación nativa porque el binario preparado en
Ubuntu~22 depende de GLIBC~2.35. La comprobación posterior de cada instalación
es de solo lectura. Clippy aprobó el workspace y todos los targets con
advertencias denegadas. Los gates globales y la integración permanecen
pendientes. Los resultados no acreditan detección de malware, autenticidad
jurídica ni ampliación de la política de soportes PDF/DOCX. El contrato y la
distinción entre evidencia local y cierre global se conservan en
\rutarepo{docs/document-upload-admission-api.md} y en
\rutarepo{docs/adr/0058-bounded-general-document-admission.md}.
